\subsection{Upper integral of a lower semi-continuous positive function}

Let X be a locally compact space, \(\mu\) a positive measure on X; we know that \(\mu\) is an increasing function on the lattice \(\mathcal{K}_{+}(\mathrm{X})\) (which will also be denoted \(K_{+}\) ).

We denote by \(\mathcal{I}_{+}(\mathrm{X})\) (or simply \(I_{+}\) ) the set of numerical functions on X, finite or not, that are positive and lower semi-continuous on X. \(^{1}\) Recall that the sum of any family of functions in \(I_{+}\) belongs to \(I_{+}\) ; the product of a function in \(I_{+}\) by a finite number \(\alpha > 0\) belongs to \(I_{+}\) ; the upper envelope of any family of functions in \(I_{+}\) and the lower envelope of a finite family of functions in \(I_{+}\) also belong to \(I_{+}\) (GT, IV, §6, No.~2, Prop. 2 and Th. 4). We shall also make use of the following lemma:

Lemma. --- Every function \(f \in I_{+}\) is the upper envelope of the set (directed for the relation \(\leqslant\) ) of all functions \(g \in K_{+}\) such that \(g \leqslant f\) .

For every \(x \in X\) such that \(f(x) > 0\) , and for every real number a such that 0 < a < \(f(x)\) , there exists by hypothesis a compact neighborhood V of x such that \(f(y) \geqslant a\) on V; on the other hand, there exists a function \(g \in K_{+}\) , with support contained in V, equal to a at the point x and \(\leqslant a\) on V (GT, IX, §1, No.~5, Th. 2); therefore \(0 \leqslant g \leqslant f\) and \(g(x) \geqslant a\) , which proves the lemma.

DEFINITION 1. --- Given a positive measure \(\mu\) on X, one calls upper integral of a function \(f \in \mathcal{I}_{+}\) (with respect to \(\mu\) ) the positive number (finite or equal to \(+\infty\) )

\[
\mu^ {*} (f) = \sup _ {g \in \mathcal {K} _ {+}, g \leqslant f} \mu (g).
\]

For every function \(f \in \mathcal{K}_+\) , it is clear that \(\mu^*(f) = \mu(f)\) , in other words \(\mu^*\) is an extension of \(\mu\) to \(\mathcal{I}_+\) .

Example. --- Let X be a discrete space, \(\mu\) a positive measure on X, and set \(\alpha(x)=\mu(\varphi_{\{x\}})\) for all \(x\in X\) . Every numerical function f defined on X is then continuous; for such a function \(f\geqslant0\) , \(\mu^{*}(f)=\sum_{x\in X}\alpha(x)f(x)\) , where we agree to set \(\alpha(x)f(x)=0\) whenever \(\alpha(x)=0\) and \(f(x)=+\infty\) . For, \(\sum_{x\in X}\alpha(x)f(x)=\sup_{M}\left(\sum_{x\in M}\alpha(x)f(x)\right)\) , where M runs over the set of all finite subsets of X. If there exists an \(x_{0}\in X\) such that \(f(x_{0})=+\infty\) and \(\alpha(x_{0})>0\) , then \(\sum_{x\in M}\alpha(x)f(x)=+\infty\) whenever \(x_{0}\in M\) , and, on the other hand, \(f\geqslant n\cdot\varphi_{\{x_{0}\}}\) for every integer n\textgreater0, therefore \(\mu^{*}(f)\geqslant n\alpha(x_{0})\) and so \(\mu^{*}(f)=+\infty\) . If, on the contrary, \(\alpha(x)=0\) at all the points where \(f(x)=+\infty\) , then the function g equal to f at the points \(x\in M\) where \(\alpha(x)>0\) and to 0 elsewhere belongs to \(H_{+}\) , and so, by virtue of the conventions made, \(\mu(g)=\sum_{x\in M}\alpha(x)f(x)\) , which again proves the relation \(\mu^{*}(f)=\sum_{x\in X}\alpha(x)f(x)\) .

PROPOSITION 1. --- For every finite real number \(\alpha > 0\) and every function \(f \in \mathcal{I}_{+}\) ,

\[
\mu^ {*} (\alpha f) = \alpha \mu^ {*} (f).\tag{1}
\]

PROPOSITION 2. --- On the set \(\mathcal{I}_{+}\) , the function \(\mu^{*}\) is increasing.

The proofs are immediate from Def. 1.

THEOREM 1. --- Let H be a nonempty set of functions in \(\mathcal{I}_{+}\) , directed for the relation \(\leqslant\) . For every positive measure \(\mu\) on X,

\[
\mu^ {*} \left(\sup _ {g \in \mathrm{H}} g\right) = \sup _ {g \in \mathrm{H}} \mu^ {*} (g) = \lim _ {g \in \mathrm{H}} \mu^ {*} (g).\tag{2}
\]

Let \(f = \sup_{g \in H} g\) . We shall first prove the theorem for the special case that the functions \(g \in H\) and their upper envelope \(f\) belong to \(\mathcal{H}_+\) . It then follows from Dini's theorem (GT, X, §4, No.~1, Th. 1) that the section filter of H converges uniformly to \(f\) on every compact subset of X, and in particular on the support K of f.~Since \(0 \leqslant g \leqslant f\) for every function \(g \in H\) , the support of every function in H is contained in K; but by definition \(\mu\) is continuous on the vector space \(\mathcal{K}(\mathrm{X},\mathrm{K};\mathbf{C})\) of continuous functions with support contained in K, for the topology of uniform convergence; whence the relation (2) in this case.

Let us pass to the general case. It is clear that \(\mu^{*}(g) \leqslant \mu^{*}(f)\) for every function \(g \in H\) . By Def. 1, it all comes down to showing that, for every function \(\psi \in K_{+}\) such that \(\psi \leqslant f\) ,

\[
\mu (\psi) \leqslant \sup _ {g \in \mathrm{H}} \mu^ {*} (g).
\]

For every function \(g \in \mathbf{H}\) , let \(\Phi_g\) be the set of functions \(\varphi \in \mathcal{K}_+\) such that \(\varphi \leqslant g\) , and let \(\Phi\) be the union of the sets \(\Phi_g\) as \(g\) runs over \(\mathbf{H}\) ; since \(\mathbf{H}\) is directed, so is \(\Phi\) , and \(f = \sup_{\varphi \in \Phi} \varphi\) . Since \(\psi \leqslant f\) , \(\psi\) is the upper envelope of the set of functions \(\inf(\psi, \varphi)\) as \(\varphi\) runs over \(\Phi\) ; but since \(\psi\) and the functions \(\inf(\psi, \varphi)\) belong to \(\mathcal{K}_+\) , the first part of the proof shows that \(\mu(\psi) = \sup_{\varphi \in \Phi} \mu(\inf(\psi, \varphi))\) . Now, each \(\varphi \in \Phi\) belongs to a set \(\Phi_g\) , therefore

\[
\mu \big (\inf (\psi , \varphi) \big) \leqslant \mu (\varphi) \leqslant \mu^ {*} (g) \leqslant \sup _ {g \in \mathrm{H}} \mu^ {*} (g),
\]

from which it follows at once that \(\mu(\psi) \leqslant \sup_{g \in \mathrm{H}} \mu^{*}(g)\) . We have thus proved that \(\mu^{*}(f) = \sup_{g \in \mathrm{H}} \mu^{*}(g)\) ; the relation \(\mu^{*}(f) = \lim_{g \in \mathrm{H}} \mu^{*}(g)\) is then a consequence of the monotone limit theorem (GT, IV, §5, No.~2, Th. 2).

THEOREM 2. --- If \(f_1\) and \(f_2\) are two functions in \(\mathcal{I}_+\) then

\[
\mu^ {*} (f _ {1} + f _ {2}) = \mu^ {*} (f _ {1}) + \mu^ {*} (f _ {2}).\tag{3}
\]

As \(\varphi_{1}\) (resp. \(\varphi_{2}\) ) runs over the set of functions in \(\mathcal{H}_{+}\) such that \(\varphi_{1} \leqslant f_{1}\) (resp. \(\varphi_{2} \leqslant f_{2}\) ), the functions \(\varphi_{1} + \varphi_{2}\) form a directed set (for \(\leqslant\) ) whose upper envelope is \(f_{1} + f_{2}\) . Therefore, by Th. 1,

\[
\mu^ {*} (f _ {1} + f _ {2}) = \sup \mu (\varphi_ {1} + \varphi_ {2}) = \sup \left(\mu (\varphi_ {1}) + \mu (\varphi_ {2})\right),
\]

where \((\varphi_{1},\varphi_{2})\) runs over the set of pairs of functions in \(\mathcal{H}_{+}\) such that \(\varphi_1\leqslant f_1\) and \(\varphi_{2}\leqslant f_{2}\) ; since

\[
\sup \left(\mu (\varphi_ {1}) + \mu (\varphi_ {2})\right) = \sup \mu (\varphi_ {1}) + \sup \mu (\varphi_ {2})
\]

(GT, IV, §5, No.~7, Cor. 2 of Prop. 12), the theorem is proved.

PROPOSITION 3. --- For every family \((f_{\iota})_{\iota \in \mathrm{I}}\) of functions in \(\mathcal{I}_{+}\) ,

\[
\mu^ {*} \left(\sum_ {\iota \in I} f _ {\iota}\right) = \sum_ {\iota \in I} \mu^ {*} (f _ {\iota}).\tag{4}
\]

For every finite subset J of I, it follows from Th. 2 (by induction on the number of elements of J) that \(\mu^{*}\left(\sum_{\iota \in \mathrm{J}}f_{\iota}\right) = \sum_{\iota \in \mathrm{J}}\mu^{*}(f_{\iota})\) ; as J runs over the set of finite subsets of I, the functions \(g_{\mathrm{J}} = \sum_{\iota \in \mathrm{J}}f_{\iota}\) belong to \(\mathcal{I}_{+}\) and form a directed set for the relation \(\leqslant\) , whose upper envelope is the function \(\sum_{\iota \in \mathrm{I}}f_{\iota}\) ; the proposition therefore follows from Th. 1.

PROPOSITION 4. --- Let \(f\) be a function in \(\mathcal{I}_{+}\) . The mapping \(\mu \mapsto \mu^{*}(f)\) of the set \(\mathcal{M}_{+}(\mathrm{X})\) of positive measures on \(\mathrm{X}\) , into the extended real line \(\overline{\mathbf{R}}\) , is lower semi-continuous for the vague topology on \(\mathcal{M}_{+}(\mathrm{X})\) (Ch. III, §1, No.~9).

For, this mapping is by definition the upper envelope of the mappings \(\mu \mapsto \mu(g)\) , where \(g\) runs over the set of functions in \(\mathcal{K}_{+}\) such that \(g \leqslant f\) ; and by definition of the vague topology, the mappings \(\mu \mapsto \mu(g)\) are continuous on \(\mathcal{M}(\mathrm{X})\) .

